\begin{frame}{그렇다면 오늘 50분은 무엇을 하는가}
  Isaac Sim을 못 돌리더라도, ``Isaac Sim이 빠르다''고 말할 때
  \textbf{그 빠름이 수학적으로 어디서 오는지}는 CPU 위에서,
  numpy 한 줄로 직접 재볼 수 있다.
  \vskip0.6em
  \small
  \begin{tabular}{@{}cl@{}}
    \toprule
    단계 & 내용 \\
    \midrule
    (1) & CPU가 환경 스텝을 \textbf{``하나씩''}만 돌릴 수 있는
      이유를 \textbf{숫자}로 확인 \\
    (2) & Isaac Sim이 그 한계를 어떻게 우회하는지
      (모든 상태를 한 배열에 담아 \textbf{커널을 한 번} 호출)
      구조도로 보고 \\
    (3) & 그 아이디어를 14.1의 진자를 \textbf{1,000개 복제}해
      numpy로 직접 시뮬레이션 · 실험 \\
    \bottomrule
  \end{tabular}
  \vskip0.5em
  \begin{itemize}
    \item 실측 머신: 학교 서버 CPU (2026-08, 워밍업 200스텝 제외)
  \end{itemize}
\end{frame}
